\chapter{Causal Structure and Time}
\label{sec:causal}
% ===================================================================

\section{Closed and Open Synchronization Trees}

The synchronization tree of a term comes in two varieties,
reflecting two different notions of how a term can evolve:
autonomously, and under interaction with an environment.
The distinction is fundamental and ramifies throughout the framework.

\begin{definition}[Closed Synchronization Tree]
The \emph{closed synchronization tree} $\mathcal{ST}_c(P)$ of a
term $P \in \terms(\GSLT)$ is the directed graph whose:
\begin{itemize}[itemsep=2pt]
  \item nodes are equivalence classes of terms reachable from $P$
        by autonomous rewrites $\rewrite^*$, requiring no external
        context;
  \item edges are labeled by rewrite rule instances: there is an
        edge $Q \xrightarrow{r} Q'$ whenever $Q \rewrite_r Q'$
        by firing rule $r$ without any surrounding context.
\end{itemize}
A term has a nontrivial closed synchronization tree if and only if
it contains a \emph{redex} --- a sub-term that matches the
left-hand side of some rule $r$ without context.
Terms with no available autonomous rewrite (normal forms, or terms
awaiting a communication partner) have a trivial closed tree
consisting of a single node.
\end{definition}

\begin{definition}[Open Synchronization Tree]
The \emph{open synchronization tree} $\mathcal{ST}_o(P)$ of a
term $P \in \terms(\GSLT)$ is the directed graph whose:
\begin{itemize}[itemsep=2pt]
  \item nodes are equivalence classes of terms reachable from $P$
        via transitions in the Milner--Sewell--Leifer labeled
        transition system;
  \item edges are labeled by minimal contexts: there is an edge
        $Q \xrightarrow{K} Q'$ whenever $K[Q] \rewrite Q'$ for
        minimal context $K$.
\end{itemize}
Every term has a potentially rich open synchronization tree,
even if its closed tree is trivial: a term that does nothing
alone may exhibit complex behavior when placed in the right
environment.
\end{definition}

\begin{remark}[The Program/Environment Boundary]
The closed tree records what a term does as a \emph{program} ---
its autonomous computational behavior.
The open tree records what a term does as a function of its
\emph{environment} --- the full range of behaviors it can exhibit
when placed in contexts.
A term with a trivial closed tree but rich open tree is a pure
interface: it does nothing alone but responds to everything.
In the Rho calculus, $x!(Q)$ is exactly this: it has no
autonomous rewrite, but in the context
$\mathtt{for}(y \leftarrow x)P \mid [-]$ it fires immediately.
\end{remark}

\begin{remark}[Relation to the LTS]
The open synchronization tree is precisely the tree unfolding of
the Milner--Sewell--Leifer labeled transition system.
The minimal contexts labeling its edges are the same minimal
contexts that generate the context-decorated HML of
Section~\ref{sec:lts-hml}.
The open tree and the HML logic are two faces of the same
structure: the tree is the extensional record of all transitions;
the logic is the intensional language for describing them.
\end{remark}

\section{The Synchronization Trees as Causal Graphs}

\begin{definition}[Synchronization Trees as Causal Graphs]
Both $\mathcal{ST}_c(P)$ and $\mathcal{ST}_o(P)$ are directed
graphs whose transitive closures are \emph{partial orders}.
Each is a \emph{causal set} in the sense of Bombelli--Lee--Sorkin:
a locally finite partial order in which $Q \leq Q'$ means
``$Q$ can causally influence $Q'$.''
The closed tree encodes autonomous causal influence;
the open tree encodes interactive causal influence mediated
by the environment.
\end{definition}

\begin{proposition}
The transitive closure of $\mathcal{ST}_c(P)$ is a partial order
on autonomously reachable terms.
The transitive closure of $\mathcal{ST}_o(P)$ is a partial order
on interactively reachable terms and extends the closed order:
$\mathcal{ST}_c(P) \subseteq \mathcal{ST}_o(P)$ as directed graphs.
\end{proposition}

\section{Causal Time}

\begin{definition}[Causal Distance]
The \emph{autonomous causal distances} are path lengths in
$\mathcal{ST}_c(P)$:
\begin{align}
  d^c_{\min}(P, Q) &= \min\bigl\{|\gamma| \;\mid\;
    \gamma \text{ is an autonomous rewrite path from } P \text{ to } Q\bigr\} \\
  d^c_{\max}(P, Q) &= \max\bigl\{|\gamma| \;\mid\;
    \gamma \text{ is an autonomous rewrite path from } P \text{ to } Q\bigr\}
\end{align}
The \emph{interactive causal distances} are path lengths in
$\mathcal{ST}_o(P)$:
\begin{align}
  d^o_{\min}(P, Q) &= \min\bigl\{|\gamma| \;\mid\;
    \gamma \text{ is a context-labeled path from } P \text{ to } Q\bigr\} \\
  d^o_{\max}(P, Q) &= \max\bigl\{|\gamma| \;\mid\;
    \gamma \text{ is a context-labeled path from } P \text{ to } Q\bigr\}
\end{align}
where $|\gamma|$ is the number of steps in $\gamma$, equal to the
nesting depth of the context labels along the path.
\end{definition}

\begin{remark}
$d^c_{\min}$ is \emph{autonomous proper time}: the irreducible
causal depth an agent traverses without environmental assistance.
$d^o_{\min}$ is \emph{interactive proper time}: the minimum number
of environmental interactions required to reach a target term.
The width of an interval in either order measures causal
concurrency, directly related to the width of parallel composition
in the Rho calculus.
The two notions coincide for terms where all behavior is eventually
triggered autonomously, and diverge for terms that are pure
interfaces awaiting environmental interaction.
\end{remark}

% ===================================================================
